% ---------------------------------------------------------------- Theory (P2)
\section{Optimal stopping and American options}

\paragraph{The problem.}
Let $(S_j)_{j=0}^{H}$ be a price process observed at dates $t_0<\dots<t_H$, with
filtration $(\mathcal F_j)$ generated by its history. Stopping at date $j$ pays a reward
$g_j$ that is known at date $j$; for an American (Bermudan) put under the risk-neutral
measure $\mathbb Q$, $g_j=e^{-rt_j}(K-S_j)^+$. A \emph{stopping time} is a random date
$\tau$ with $\{\tau=j\}\in\mathcal F_j$: the decision to stop uses only the past. The price is
\begin{equation}
V_0=\sup_{\tau\in\mathcal T_{0,H}}\ \mathbb E^{\mathbb Q}\big[g_\tau\big].
\end{equation}

\paragraph{Snell envelope.}
Define backwards
\begin{equation}
U_H=g_H,\qquad U_j=\max\big(g_j,\ \mathbb E[U_{j+1}\mid\mathcal F_j]\big),\quad j=H-1,\dots,0 .
\label{eq:snell}
\end{equation}
The process $(U_j)$ is the smallest supermartingale that dominates $(g_j)$; the stopping
time $\tau^*=\min\{j:U_j=g_j\}$ is optimal and $V_0=U_0$~\cite{shreve2}. The term
$C_j=\mathbb E[U_{j+1}\mid\mathcal F_j]$ is the \emph{continuation value}: one exercises
as soon as the immediate reward is at least the value of waiting.

\paragraph{The Markov case.}
If $S$ is Markov and $g_j=G_j(S_j)$, the continuation value is a function of today's state
only, $C_j=c_j(S_j)$. The binomial tree computes $c_j$ exactly on a lattice; Longstaff and
Schwartz estimate it by regressing realised discounted cash flows on functions of
$S_j$~\cite{ls2001}; Deep Optimal Stopping learns the exercise decision at each date with a
neural network of $S_j$~\cite{becker2019}. When the dynamics have memory (as under rough
volatility), $C_j$ depends on the whole path $(S_0,\dots,S_j)$ and these methods lose
information.

\paragraph{Continuous time: a free-boundary problem.}
When exercise is allowed at all times, a put is exercised when $S_t$ falls below a boundary
$L(t)$. Above it the value solves the Black--Scholes equation
$\partial_tV+rS\partial_SV+\tfrac12\sigma^2S^2\partial_{SS}V-rV=0$; below it $V=K-S$; and at
$S=L(t)$ the value meets the payoff smoothly, $\partial_SV=-1$ (smooth pasting). The
boundary is unknown in advance. For a perpetual put it is constant~\cite{shreve2},
\begin{equation}
L^*=\frac{2rK}{2r+\sigma^2},\qquad V(S)=(K-L^*)\Big(\frac{S}{L^*}\Big)^{-2r/\sigma^2}\ \ (S>L^*);
\end{equation}
with $K=40$, $r=0.06$, $\sigma=0.2$ this gives $L^*=30$ and $V(36)=5.79$.

\paragraph{Solving the finite-maturity problem numerically.}
For a finite maturity there is no closed form, and the free-boundary problem is written as
a linear complementarity problem: with $\mathcal L V=(r-\delta)S\partial_SV+\tfrac12\sigma^2S^2\partial_{SS}V-rV$,
\begin{equation}
\min\big(-\partial_tV-\mathcal LV,\ V-g\big)=0,\qquad V(T,S)=g(S).
\end{equation}
Where waiting is optimal the first term vanishes, which is the Feynman--Kac equation for
the discounted expected payoff; where exercise is optimal $V=g$. We solve it in
$x=\log S$, where the coefficients are constant, with Crank--Nicolson time steps (the first
four fully implicit, to damp the oscillations caused by the kink of the payoff) and impose
$V\ge g$ by a penalty iteration~\cite{forsyth2002}: each tridiagonal system is solved with
a large penalty on the nodes where $V<g$, until those nodes stop changing. A Bermudan
option is priced by solving the European equation between exercise dates and taking
$\max(V,g)$ on each of them. Table~\ref{tab:benchmarks} compares the result with a
5{,}000-step-per-year binomial tree: the two independent methods agree to within $0.0004$
on all twenty cases, so either can serve as the ground truth for the simulation methods.
The early-exercise premium is the gap between the American and European columns, and the
Bermudan price with 50 exercise dates per year is only about $0.007$ below the American
one. For $S_0=36$, $\sigma=0.2$, $T=1$ the computed boundary rises from $L(0)=32.9$ to
$L(T^-)=K=40$ and stays above the perpetual boundary $L^*=30$ throughout: with less time
left, waiting is worth less, so exercise starts at higher prices.

\begin{table}[ht]
\centering\small
\caption{American put, $K=40$, $r=0.06$: European (Black--Scholes), Bermudan (50 exercise
dates per year) and American prices from the binomial tree and the finite-difference
solver of the free-boundary problem.}
\label{tab:benchmarks}
\begin{tabular}{rrr r rr rr}
\toprule
$S_0$ & $\sigma$ & $T$ & European & \multicolumn{2}{c}{Bermudan} & \multicolumn{2}{c}{American}\\
 & & & & tree & PDE & tree & PDE\\
\midrule
% S0 & sigma & T & European & Bermudan tree & Bermudan PDE & American tree & American PDE   (generated by experiments/benchmarks.py)
36 & 0.2 & 1 & 3.844 & 4.478 & 4.478 & 4.487 & 4.487 \\
36 & 0.2 & 2 & 3.763 & 4.840 & 4.840 & 4.848 & 4.848 \\
36 & 0.4 & 1 & 6.711 & 7.101 & 7.101 & 7.109 & 7.109 \\
36 & 0.4 & 2 & 7.700 & 8.507 & 8.507 & 8.514 & 8.514 \\
38 & 0.2 & 1 & 2.852 & 3.250 & 3.250 & 3.257 & 3.257 \\
38 & 0.2 & 2 & 2.991 & 3.745 & 3.745 & 3.751 & 3.751 \\
38 & 0.4 & 1 & 5.834 & 6.148 & 6.148 & 6.155 & 6.154 \\
38 & 0.4 & 2 & 6.979 & 7.668 & 7.668 & 7.675 & 7.675 \\
40 & 0.2 & 1 & 2.066 & 2.314 & 2.314 & 2.320 & 2.319 \\
40 & 0.2 & 2 & 2.356 & 2.885 & 2.884 & 2.890 & 2.890 \\
40 & 0.4 & 1 & 5.060 & 5.312 & 5.312 & 5.318 & 5.318 \\
40 & 0.4 & 2 & 6.326 & 6.917 & 6.917 & 6.923 & 6.923 \\
42 & 0.2 & 1 & 1.465 & 1.617 & 1.617 & 1.621 & 1.621 \\
42 & 0.2 & 2 & 1.841 & 2.212 & 2.212 & 2.217 & 2.217 \\
42 & 0.4 & 1 & 4.379 & 4.583 & 4.582 & 4.588 & 4.588 \\
42 & 0.4 & 2 & 5.736 & 6.244 & 6.244 & 6.250 & 6.250 \\
44 & 0.2 & 1 & 1.017 & 1.110 & 1.110 & 1.113 & 1.113 \\
44 & 0.2 & 2 & 1.429 & 1.690 & 1.690 & 1.693 & 1.693 \\
44 & 0.4 & 1 & 3.783 & 3.948 & 3.948 & 3.953 & 3.953 \\
44 & 0.4 & 2 & 5.202 & 5.641 & 5.641 & 5.647 & 5.647 \\

\bottomrule
\end{tabular}
\end{table}

\paragraph{Optional stopping.}
If $(M_j)$ is a martingale and $\tau$ a bounded stopping time, then
$\mathbb E[M_\tau]=M_0$: no stopping rule can improve on a fair game. Two consequences
guide our experiments: stopping a standard Brownian motion is worth exactly~$0$, so memory
can only help when increments are dependent; and if market prices are close to
martingales, no exit rule can beat holding to the horizon in mean, so any benefit of
memory should appear mainly in risk-adjusted terms.
